% Zero-Shot Conservation Critics for Generative Physical Simulation
% arXiv preprint.  Build: pdflatex main && pdflatex main
%
% NOTE ON VERIFICATION: no LaTeX toolchain was available in the authoring
% environment, so this file has NOT been compiled.  Its structure is checked
% mechanically by `docs/paper/check_tex.py` (brace / environment balance,
% preamble presence, undefined-macro scan); the rendered output is unverified.
\documentclass[11pt]{article}

\usepackage[margin=1in]{geometry}
\usepackage{amsmath,amssymb}
\usepackage{booktabs}
\usepackage{graphicx}
\usepackage{microtype}
\usepackage[hidelinks]{hyperref}

\newcommand{\argmax}{\operatorname*{arg\,max}}
\newcommand{\rmse}{\mathrm{RMSE}}
\newcommand{\eps}{\varepsilon}

\title{Zero-Shot Conservation Critics\\ for Generative Physical Simulation}
\author{}
\date{}

\begin{document}
\maketitle

\begin{abstract}
Generative models of physical scenes are typically validated with a \emph{critic}: a
score that ranks or filters candidate trajectories. Learned critics are the default
choice, but they require labelled data and can be brittle. We ask whether a
\textbf{zero-shot critic built from conservation laws} --- no training, no labels ---
is competitive.

We study \emph{candidate selection}: given $K$ trajectories generated from the same
initial condition by a learned dynamics model, pick the one closest to the true
solution. This is the operation a critic performs inside a generation loop, and unlike
classification-style metrics it needs no threshold on ``how wrong is wrong''.

Across three systems --- a 6-body 3D family, the figure-eight three-body choreography,
and a 2D Kepler two-body system --- a zero-shot conservation critic attains
\textbf{100\% within-group selection accuracy} on the standard benchmarks, reducing
trajectory error by 75--82\% relative to random selection and landing within
$1.03$--$1.05\times$ of the best achievable selection. On the 6-body system (5 seeds,
95\% CI) it reaches $82.8\%\pm3.3\%$ in distribution and $84.7\%\pm1.8\%$ out of
distribution, versus $71.2\%\pm3.2\%$ and $72.2\%\pm2.5\%$ for a supervised
pairwise-ranking critic trained on the same 960 labelled candidates.

We also establish what the conservation critic is \emph{not}. It does \textbf{not} beat
a correctly-used Hamiltonian Neural Network critic: the two tie exactly (identical
selection error on both standard benchmarks). And its lead does not survive a change of
\emph{generator}: with an HNN as the generator the supervised critic is at least as good.
Claims that the conservation critic is more robust out of distribution, or that combining
it with a learned critic helps, are \textbf{falsified} here and retracted.

Three methodological findings, each independently useful. (i) \textbf{Pooled rank
correlation is misleading for critics}: a supervised regressor scored pooled Spearman
$\rho=-0.67$, comparable to the best critic, while achieving a within-group win rate of
$53\%$, i.e.\ chance. (ii) \textbf{A conservation channel's discriminative power is set
by its noise floor}, determined by the integrator and force structure rather than by the
physical importance of the quantity: momentum and angular momentum are conserved
\emph{exactly} by any pairwise central-force integrator (floor $\sim10^{-16}$) and
separate true from generated trajectories by $\sim10^{12}$, whereas energy is conserved
only to the integrator's truncation error ($O(\Delta t)$, floor $0.60$ at
$\Delta t=0.01$) and separates by only $1.13\times$. (iii) \textbf{Synthetic failure
modes must be validated before use}: a taxonomy of hand-built corruptions is not
structurally comparable to a real generator's failures --- two of five cannot be tuned
to the generator's error magnitude at all, and conservation fingerprints differ by up
to $238$ normalized units.
\end{abstract}

\section{Introduction}

A generative model of physical scenes --- a learned surrogate, a neural ODE, a video
model --- produces trajectories that may look plausible while violating physics. The
standard remedy is a \emph{critic}: a function that scores a candidate so that bad
samples can be filtered, ranked, or used as a training signal.

Learned critics dominate, and they work: a network can capture subtle dynamical errors
that no simple statistic catches. But they need labelled examples, and labels for
physical validity are exactly what one is trying to avoid needing.

This paper asks a narrower and more testable question than ``can physics replace
learning?'':

\begin{quote}
\emph{Is a critic built from conservation laws --- which is training-free and
label-free --- competitive with the strongest learned critics at selecting good
candidates?}
\end{quote}

The answer we establish is: \textbf{yes, it ties the strongest}, and it does so at zero
labelling cost. We also document three claims that turned out to be false, because the
process of falsifying them produced the paper's methodological contributions.

\paragraph{Contributions.}
\begin{itemize}
  \item \textbf{C1. A threshold-free evaluation protocol} for critics in a generation
    loop: \emph{within-group selection}. Given $K$ candidates from one initial
    condition, measure the critic's rank correlation \emph{within that group} and its
    selection error relative to random and oracle. Section~\ref{sec:metric} shows that
    the more obvious pooled metric is badly misleading.
  \item \textbf{C2. A zero-shot conservation critic that ties the strongest
    baselines}, including a correctly-used Hamiltonian Neural Network and a supervised
    pairwise-ranking critic trained on up to 4800 labelled candidates
    (Section~\ref{sec:main}).
  \item \textbf{C3. A noise-floor theory of channel selection} that explains,
    predicts, and improves the channel set (Sections~\ref{sec:nf} and
    \ref{sec:ablation}).
  \item \textbf{C4. A negative result on synthetic failure taxonomies}, with a
    calibration procedure to test comparability before the taxonomy is used
    (Section~\ref{sec:taxonomy}).
  \item \textbf{C5. Three retractions} of claims made in earlier drafts of this work
    (Section~\ref{sec:retractions}).
\end{itemize}

\section{Related work}

\paragraph{Learned physical simulators.}
Hamiltonian Neural Networks \cite{greydanus2019hnn} and Lagrangian Neural Networks
\cite{cranmer2020lNN} learn a scalar potential whose derivatives reproduce the
dynamics. They are the natural learned alternative to a hand-written conservation law,
and we adopt an HNN as our primary literature baseline.

\paragraph{Neural surrogates and operators} learn the flow map directly. Their errors are
typically evaluated by trajectory RMSE against a reference integration, which is
available in simulation but not in the field.

\paragraph{Physics-based validation.} Conservation checks are standard in numerical
relativity and molecular dynamics as diagnostics. Using them as a \emph{critic inside a
generative loop} --- the object of this paper --- is less studied.

\paragraph{Positioning.} We do not propose a new generator. We characterise when a
training-free physical critic is sufficient, and we show that the answer depends on a
numerical property (the noise floor) rather than on the physics.

\section{Setup}
\label{sec:setup}

\subsection{Task: within-group candidate selection}

For each initial condition (an ``IC''):

\begin{enumerate}
  \item integrate the true equations of motion to obtain a reference trajectory
    $x^{\star}(t)$;
  \item run a learned generator $K=8$ times from the \emph{same} IC, with varying
    sampling noise $\sigma\in\{0,\,0.02,\,0.05,\,0.10,\,0.20,\,0.35,\,0.6,\,1.0\}$;
  \item compute each candidate's error $e_j = \rmse(x_j, x^{\star})$;
  \item score all candidates with each critic and select $\argmax_j \mathrm{score}_j$.
\end{enumerate}

\begin{table}[t]
\centering
\caption{Metrics. Lower is better except for the win rate.}
\label{tab:metrics}
\begin{tabular}{llcc}
\toprule
Metric & Definition & Chance & Perfect \\
\midrule
within-group $\rho$ & mean over ICs of Spearman(scores, $e$) inside the IC & $0$ & $-1$ \\
IC win rate & fraction of ICs where the selection beats that IC's median error & 50\% & 100\% \\
selected / random & median selected error $\div$ median of per-IC mean error & 1.00 & --- \\
selected / oracle & median selected error $\div$ median of per-IC minimum error & --- & 1.00 \\
\bottomrule
\end{tabular}
\end{table}

\subsection{Systems}

\begin{table}[t]
\centering
\caption{Systems used throughout.}
\label{tab:systems}
\begin{tabular}{llll}
\toprule
System & $N$ & Dim. & Notes \\
\midrule
6-body, eccentric family & 6 & 3D & $e\sim U(0,0.3)$ ID; $e\sim U(0.5,0.8)$ OOD \\
figure-eight choreography & 3 & 2D & \cite{chenciner2000}; $T=6.32591398$ \\
Kepler two-body & 2 & 2D & standard low-dimensional HNN/LNN testbed \\
\bottomrule
\end{tabular}
\end{table}

Checker for the figure-eight initial data: after exactly one standard period at
$\Delta t = 0.002$ (3163 steps) the configuration returns to the initial condition to
within $3.29\times10^{-3}$ in position. All integrations use symplectic Euler at
$\Delta t = 0.002$--$0.01$.

\subsection{Generators}

The generator is a pairwise force model: a shared MLP applied to
$[\,r_{ij},\,|r_{ij}|,\,v_{ij},\,m_i,\,m_j\,]$ producing a scalar pair force, summed
over neighbours. It is deliberately trained on a \emph{limited budget} so that its
failures are large enough to be selectable among (one-step acceleration error $0.23$
for the 6-body system; three budgets, $0.23/0.0068/0.81$, are used for robustness
checks).

Negative samples therefore come from a \textbf{real learned generator}, not from
hand-built corruptions. Section~\ref{sec:taxonomy} explains why that choice was forced.

\subsection{Critics}

\begin{table}[t]
\centering
\caption{Critics compared. Training budgets differ per experiment; each table states
its own. ``HNN, energy drift'' is the naive way to use an HNN as a critic; see
Section~\ref{sec:main}.}
\label{tab:critics}
\begin{tabular}{lll}
\toprule
Critic & Labels required & Notes \\
\midrule
\textbf{Conservation (zero-shot)} & \textbf{0} & angular $+$ momentum channels \\
HNN, equation residual & true trajectories & \cite{greydanus2019hnn} \\
HNN, energy drift & true trajectories & naive use as a critic \\
Pairwise ranking loss & 960--4800 candidates & supervised, target-aligned \\
Absolute binary classifier & 960--4800 candidates & supervised, naive target \\
Absolute error regressor & 960--4800 candidates & supervised, naive target \\
Random / oracle & --- & reference points \\
\bottomrule
\end{tabular}
\end{table}

\section{Method: the conservation critic}

\subsection{Conservation drift as an error functional}

Let the true acceleration field be a sum of pairwise central forces,
\begin{equation}
a^{\text{true}}_i(r)=\frac{1}{m_i}\sum_{j\neq i} f(r_{ij})\,\hat{r}_{ij},
\end{equation}
and let $\delta a(t) = a^{\text{gen}}(\hat{z}(t)) - a^{\text{true}}(\hat{r}(t))$ be the
generator's acceleration error along its own trajectory $\hat{z}$. Because
$\sum_i m_i a^{\text{true}}_i \equiv 0$ and
$\sum_i m_i (r_i \times a^{\text{true}}_i) \equiv 0$ hold pointwise,
\begin{equation}
\label{eq:noether}
\frac{dP}{dt}=\sum_i m_i\,\delta a_i,\qquad
\frac{dL}{dt}=\sum_i m_i\,(\hat{r}_i\times\delta a_i),\qquad
\frac{dE}{dt}=\sum_i m_i\,(\hat{v}_i\cdot\delta a_i).
\end{equation}
The conservation residual is therefore the \textbf{projection of the accumulated force
error onto the symmetry generators} --- a physics-meaningful functional, not a
heuristic.

Semi-quantitative validation (40 ICs, medians of the ratio prediction/measured):
momentum $0.53\times$, angular $1.23\times$, energy $9.0\times$; log-correlations
$0.72/0.31/0.22$. Equation~\eqref{eq:noether} is exact; its integrated form is
validated to within a factor of a few for the channels with a clean numerical floor.

\begin{quote}
\textbf{What is \emph{not} claimed.} Equation~\eqref{eq:noether} bounds the
\emph{drift} by the force error. It gives \textbf{no bound on trajectory error}: any
$\delta a$ lying in the orthogonal complement of the three generators produces zero
residual with arbitrarily large trajectory deviation. We never assume such a
$\delta a$ is absent.
\end{quote}

\subsection{Which channels to use, and why}
\label{sec:nf}

The channels are measured on the reference trajectories. linear and angular momentum
are exact invariants of any pairwise central force, to roundoff, independent of the
integrator; their floor is $\eta_P \sim \eta_L \sim \eps_{\text{mach}} \approx
10^{-16}$. Energy is conserved only to the integrator's truncation error; for the
first-order symplectic Euler used here, $\eta_E = O(\Delta t)$.

This predicts the measured channel ranking and is confirmed by a step-size sweep
(identical generator, only $\Delta t$ varied):

\begin{table}[t]
\centering
\caption{Noise floors and energy separation as a function of step size. The energy
floor halves with $\Delta t$; the momentum and angular floors stay pinned at roundoff.}
\label{tab:noisefloor}
\begin{tabular}{rrrrr}
\toprule
$\Delta t$ & energy floor & angular floor & momentum floor & energy separation \\
\midrule
0.0100 & $6.02\mathrm{e}{-01}$ & $4.01\mathrm{e}{-16}$ & $2.88\mathrm{e}{-16}$ & 0.81 \\
0.0050 & $\mathbf{2.65\mathrm{e}{-01}}$ & $5.06\mathrm{e}{-16}$ & $3.68\mathrm{e}{-16}$ & 1.84 \\
0.0025 & $\mathbf{1.36\mathrm{e}{-01}}$ & $6.62\mathrm{e}{-16}$ & $4.39\mathrm{e}{-16}$ & \textbf{3.58} \\
\bottomrule
\end{tabular}
\end{table}

\textbf{Channel utility is a numerical property, not a physical one.} The critic
therefore uses angular $+$ momentum by default, and excludes both the energy channel
(weak floor) and the centre-of-mass channel (an exact invariant, but its variation is
nearly uncorrelated with shape error --- measured IC win rate $16.7$--$24.2\%$, below
chance).

\section{Experiments}
\label{sec:experiments}

\subsection{Why the obvious metric is wrong}
\label{sec:metric}

Reported ``critic quality'' is often a rank correlation pooled over all candidates. That
metric conflates two effects, because candidates sampled with larger $\sigma$ are worse
across the board. A critic that only detects the overall noise level scores well pooled
while being unable to rank inside a group. Measured on the 6-body system (600 training
ICs $=$ 4800 labelled candidates):

\begin{table}[t]
\centering
\caption{Pooled versus within-group rank correlation. The supervised error regressor
looks competitive pooled but has a within-group win rate of $53\%$ --- chance.}
\label{tab:pooled}
\begin{tabular}{lrrrr}
\toprule
Critic & pooled $\rho$ & \textbf{within-group $\rho$} & \textbf{IC win rate} & selected / random \\
\midrule
Conservation & $-0.7044$ & $\mathbf{-0.6881}$ & $\mathbf{84.0\%}$ & $\mathbf{0.637\times}$ \\
Pairwise ranking (supervised) & $-0.4689$ & $-0.5786$ & $74.5\%$ & $0.656\times$ \\
\textbf{Error regressor (supervised)} & $\mathbf{-0.6685}$ & $\mathbf{-0.3799}$ & $\mathbf{53.0\%}$ & $0.758\times$ \\
Logistic fusion of all three & $-0.4509$ & $-0.6825$ & $83.5\%$ & $0.641\times$ \\
\bottomrule
\end{tabular}
\end{table}

A related effect: aligning the training target with the evaluation is worth more than
model capacity. Identical architecture, two targets --- absolute error regression
reaches an IC win rate of $\mathbf{47.5\%}$; IC-relative error regression reaches
$\mathbf{86.5\%}$, a 39-point gap from the target alone.

\subsection{Main results}
\label{sec:main}

\begin{table}[t]
\centering
\caption{figure-eight three-body ($N=3$, 2D), 50 test ICs, $K=8$. ``Separation'' is the
relative gap between the critic's score on the true trajectory and on the $\sigma=0$
candidate.}
\label{tab:fig8}
\begin{tabular}{lrrrrr}
\toprule
Critic & separation & within-group $\rho$ & IC win rate & sel./random & sel./oracle \\
\midrule
\textbf{Conservation (zero-shot)} & $\mathbf{8.7\mathrm{e}{11}}$ & $-0.9486$ & $\mathbf{100.0\%}$ & $\mathbf{0.250\times}$ & $\mathbf{1.03\times}$ \\
HNN, energy drift (naive) & $0.076$ & $-0.5405$ & $50.0\%$ & $0.544\times$ & $2.24\times$ \\
\textbf{HNN, equation residual} & $0.017$ & $\mathbf{-0.9524}$ & $\mathbf{100.0\%}$ & $\mathbf{0.250\times}$ & $\mathbf{1.03\times}$ \\
\midrule
random / oracle & --- & --- & 50.0\% / 100\% & $1.000\times$ / $0.243\times$ & $4.11\times$ / $1.00\times$ \\
\bottomrule
\end{tabular}
\end{table}

\begin{table}[t]
\centering
\caption{Kepler two-body ($N=2$, 2D), 50 test ICs, $K=8$.}
\label{tab:kepler}
\begin{tabular}{lrrrrr}
\toprule
Critic & separation & within-group $\rho$ & IC win rate & sel./random & sel./oracle \\
\midrule
\textbf{Conservation (zero-shot)} & $\mathbf{2.4\mathrm{e}{14}}$ & $-0.9495$ & $\mathbf{100.0\%}$ & $\mathbf{0.185\times}$ & $\mathbf{1.05\times}$ \\
HNN, energy drift (naive) & $0.034$ & $-0.5705$ & $52.0\%$ & $0.463\times$ & $2.63\times$ \\
\textbf{HNN, equation residual} & $0.058$ & $\mathbf{-0.9490}$ & $\mathbf{100.0\%}$ & $\mathbf{0.185\times}$ & $\mathbf{1.05\times}$ \\
\midrule
random / oracle & --- & --- & 50.0\% / 100\% & $1.000\times$ / $0.176\times$ & $5.68\times$ / $1.00\times$ \\
\bottomrule
\end{tabular}
\end{table}

\textbf{The conservation critic ties the HNN exactly} on both systems --- identical
selection error, identical win rate. The learned critic requires training an HNN; the
conservation critic requires nothing.

Note also the separation column: the conservation critic separates true from generated
trajectories by $10^{11}$--$10^{14}$, the HNN by $\sim10^{-2}$. Yet the HNN still ranks
correctly inside a group, because its residual varies monotonically with $\sigma$.
Large absolute separation and good within-group ranking are different properties.

\begin{table}[t]
\centering
\caption{6-body system, 5 seeds, 95\% CI (120 training ICs). The conservation critic's
advantage over the supervised ranking critic has non-overlapping 95\% CIs on both
splits.}
\label{tab:6body}
\begin{tabular}{lrrrr}
\toprule
Critic & ID within-group $\rho$ & ID win rate & OOD win rate & training samples \\
\midrule
\textbf{Conservation (zero-shot)} & $\mathbf{-0.6404\pm0.0459}$ & $\mathbf{82.8\%\pm3.3\%}$ & $\mathbf{84.7\%\pm1.8\%}$ & \textbf{0} \\
Pairwise ranking (supervised) & $-0.5627\pm0.0471$ & $71.2\%\pm3.2\%$ & $72.2\%\pm2.5\%$ & 960 \\
Absolute binary (naive) & $-0.4958\pm0.0768$ & $62.0\%\pm8.8\%$ & $67.0\%\pm7.3\%$ & 960 \\
random & --- & 50.0\% & 50.0\% & --- \\
\bottomrule
\end{tabular}
\end{table}

The numbers in Table~\ref{tab:6body} come from 120 training ICs. In
Section~\ref{sec:metric} and Section~\ref{sec:retractions} the supervised critics were
given 600 ICs (4800 labelled candidates) --- five times more --- and the conclusion is
unchanged. We quote each experiment's own budget rather than a single figure.

\subsection{Ablation over channel sets}
\label{sec:ablation}

Four settings (two generators $\times$ ID/OOD), reported as selected/random (lower is
better):

\begin{table}[t]
\centering
\caption{Channel-set ablation. The previously-used equal-weight sum of four channels is
substantially suboptimal.}
\label{tab:ablation}
\begin{tabular}{lrrrr}
\toprule
Channel set & G\_A ID & G\_A OOD & G\_B ID & G\_B OOD \\
\midrule
angular & $0.640\times$ & $\mathbf{0.637\times}$ & $\mathbf{0.056\times}$ & $0.063\times$ \\
\textbf{angular $+$ momentum (adopted)} & $0.640\times$ & $0.653\times$ & $\mathbf{0.056\times}$ & $\mathbf{0.059\times}$ \\
four channels, equal weight (previous) & $0.629\times$ & $0.713\times$ & $0.067\times$ & $0.094\times$ \\
energy only & $0.794\times$ & $0.828\times$ & $0.308\times$ & $0.194\times$ \\
centre of mass only & $\mathbf{1.038\times}$ & $\mathbf{1.057\times}$ & $\mathbf{1.327\times}$ & $\mathbf{1.575\times}$ \\
\bottomrule
\end{tabular}
\end{table}

On the stronger generator the old four-channel sum is $1.21\times$ of oracle where
angular$+$momentum is $1.02\times$, and out of distribution the gap widens to
$1.61\times$ versus $1.02\times$. The centre-of-mass channel is worse than chance on its
own.

\subsection{A negative result: synthetic failure modes are not comparable}
\label{sec:taxonomy}

Critic rankings are often measured on hand-built corruptions (add noise, swap in a bad
integrator, freeze, splice, time-warp). We tested whether such a taxonomy is comparable
to a real generator's failures. \textbf{It is not.}

\paragraph{(a) Two of five corruptions cannot be calibrated.} Tuning severity so that
each corruption's median RMSE matches the generator's ($0.01468$):

\begin{table}[t]
\centering
\caption{Severity calibration. ``Lower bound'' means the corruption was made as mild as
its parameterisation allows.}
\label{tab:calib}
\begin{tabular}{lrrl}
\toprule
Corruption & calibrated severity & achieved RMSE & verdict \\
\midrule
noise & $0.740$ & $0.01477$ & calibrated \\
freeze & $0.428$ & $0.01342$ & calibrated \\
warp & $0.302$ & $0.01342$ & calibrated \\
\textbf{euler} & lower bound & $\mathbf{0.14127}$ & \textbf{$10\times$ too destructive} \\
\textbf{splice} & lower bound & $\mathbf{0.59365}$ & \textbf{$40\times$ too destructive} \\
\bottomrule
\end{tabular}
\end{table}

\texttt{splice} is worse than merely uncalibrated: it performs a discrete particle roll
and its severity parameter has \textbf{no effect at all}. It is an artifact of the
construction, not a failure mode.

\paragraph{(b) The conservation fingerprints have different shape.}

\begin{table}[t]
\centering
\caption{Separation of true from corrupted trajectories, per channel. The real
generator's failure mode is breaking Newton's third law, which destroys angular and
linear momentum conservation. Four of the five synthetic corruptions leave those
quantities at $0.9$--$1.1$.}
\label{tab:fingerprint}
\begin{tabular}{lrrrrrr}
\toprule
Channel & noise & euler & freeze & splice & warp & \textbf{real generator} \\
\midrule
energy & $3.89$ & $0.00$ & $0.86$ & $0.98$ & $0.93$ & $\mathbf{1.02}$ \\
angular & $7.5\mathrm{e}{15}$ & $3.1\mathrm{e}{10}$ & $0.95$ & $0.97$ & $1.13$ & $\mathbf{1.2\mathrm{e}{15}}$ \\
momentum & $6.3\mathrm{e}{15}$ & $3.9\mathrm{e}{8}$ & $0.87$ & $0.96$ & $0.94$ & $\mathbf{1.9\mathrm{e}{16}}$ \\
centre of mass & $1.10$ & $0.02$ & $0.85$ & $1.03$ & $1.03$ & $\mathbf{1.10}$ \\
\bottomrule
\end{tabular}
\end{table}

Normalized fingerprint distances range from $7.4$ to $238$.

\paragraph{Consequence.} No result in this paper is computed on a synthetic corruption
set. All negative samples come from real learned generators. We recommend the
calibration procedure above as a prerequisite for anyone using such a taxonomy.

\subsection{Cross-generator check: the ranking depends on the generator}
\label{sec:crossgen}

The results above use a pairwise force model as the generator. A reviewer's first
question is whether they are an artefact of that choice. We re-ran the comparison with a
\textbf{Hamiltonian Neural Network used as the generator} \cite{greydanus2019hnn}, which
has a fundamentally different inductive bias: it learns a black-box scalar $H(q,p)$ and
derives the vector field from Hamilton's equations, so it is \emph{not}
translation- or rotation-invariant, unlike a pairwise force model.

\begin{table}[t]
\centering
\caption{HNN used as the \emph{generator}. The conservation critic's lead does not
survive the change of generator.}
\label{tab:crossgen}
\begin{tabular}{llrrrr}
\toprule
System & Critic & within-group $\rho$ & IC win rate & sel./random & sel./oracle \\
\midrule
figure-eight & Conservation (zero-shot) & $-0.6881$ & $78.0\%$ & $0.864\times$ & $1.10\times$ \\
figure-eight & \textbf{Pairwise ranking (supervised)} & $\mathbf{-0.8890}$ & $\mathbf{96.0\%}$ & $\mathbf{0.794\times}$ & $\mathbf{1.01\times}$ \\
Kepler & Conservation (zero-shot) & $-0.8048$ & $\mathbf{98.0\%}$ & $0.574\times$ & $1.06\times$ \\
Kepler & \textbf{Pairwise ranking (supervised)} & $\mathbf{-0.9014}$ & $96.0\%$ & $\mathbf{0.552\times}$ & $\mathbf{1.02\times}$ \\
\midrule
\multicolumn{2}{l}{random (figure-eight / Kepler)} & --- & 50.0\% & $1.000\times$ & $1.27\times$ / $1.84\times$ \\
\bottomrule
\end{tabular}
\end{table}

\textbf{The conservation critic no longer leads.} On the figure-eight it is clearly
behind ($78.0\%$ versus $96.0\%$); on Kepler the two are close with the supervised critic
marginally ahead in selection error. Both remain far better than random
($0.55$--$0.86\times$ versus $1.00\times$).

The conservation fingerprint shows the HNN generator still breaks angular-momentum
conservation severely (separation $2.4\times10^{13}$ and $7.3\times10^{12}$), so the
critic's signal is not weakened.

\paragraph{A candidate explanation, tested and rejected.}
The conservation residual is a projection of the force error onto the symmetry generators
(Eq.~\ref{eq:noether}), so one might expect the \emph{full} force error to carry strictly
more information, with the HNN generator's advantage lying in structure that conservation
cannot see. We tested this directly: for every candidate we measured the accumulated force
error $\int\lVert a^{\text{gen}}-a^{\text{true}}\rVert\,dt$ along its own trajectory and
correlated it, within each IC, with the true trajectory error.

\begin{table}[t]
\centering
\caption{Tested explanation. Positive Spearman means the diagnostic predicts trajectory
error. The full force error is a \emph{worse} predictor than the conservation residual in
all four settings, so this diagnostic cannot support the ``invisible structure''
explanation.}
\label{tab:decomp}
\begin{tabular}{lrrr}
\toprule
Setting & $\rho$(conservation) & $\rho$(force error) & gap \\
\midrule
figure-eight $\times$ pairwise & $+0.9524$ & $+0.1702$ & $-0.7821$ \\
figure-eight $\times$ HNN & $+0.7768$ & $+0.1339$ & $-0.6429$ \\
Kepler $\times$ pairwise & $+0.8173$ & $+0.1238$ & $-0.6935$ \\
Kepler $\times$ HNN & $+0.7929$ & $+0.2083$ & $-0.5845$ \\
\bottomrule
\end{tabular}
\end{table}

\textbf{The explanation does not hold.} The force-error magnitude is a far weaker
predictor than the conservation residual everywhere, so nothing here supports attributing
the HNN generator's advantage to information that conservation discards. Two reasons the
diagnostic may simply be the wrong instrument, neither of which we separated: the force
error is evaluated at the \emph{generated} positions, so a candidate that has drifted is
compared in a different region of state space; and trajectory error is the \emph{second}
time integral of $\delta a$, so oscillating force errors that largely cancel leave little
trace --- a magnitude is not the right upper bound on usable information.

We therefore leave the mechanism \textbf{unexplained}. What stands is the measurement:
with an HNN generator the supervised critic is at least as good as the conservation
critic. Only the proposed explanation failed.

\paragraph{What this changes.} The claim must be stated per generator, not universally:

\begin{quote}
A zero-shot conservation critic is a \textbf{reliable but not universally optimal}
supervision signal. For pairwise-force generators it leads the strongest supervised
critic; for a Hamiltonian-parameterised generator it is matched or slightly beaten.
\end{quote}

This does not contradict the tie reported in Section~\ref{sec:main}: that compares
\emph{critics} under a fixed pairwise-force generator, whereas this compares
\emph{generators}. Together they place the conservation critic in the same band as the
strongest learned supervision, with the ordering depending on the setting rather than
following a general rule.

\section{Limitations}

\begin{enumerate}
  \item \textbf{Generators are learned surrogate force models}, not diffusion or
    autoregressive video models. The qualitative structure of the failures (breaking
    equal-and-opposite forces) is plausible for any pairwise-parameterised model, but we
    did not verify it on a large pretrained generative model.
  \item \textbf{The generator's parameterisation changes the ranking.}
    Section~\ref{sec:crossgen} shows that switching from a pairwise force model to a
    Hamiltonian Neural Network removes the conservation critic's lead. The claim must be
    stated per generator; we have two architectures, not a spectrum.
  \item \textbf{Ground truth is always available} here, because the systems are
    simulated. The selection task is therefore well-defined; in the field one normally
    lacks $x^{\star}$. Our critic does not need $x^{\star}$ at \emph{use} time --- only
    our \emph{evaluation} does.
  \item \textbf{Three systems, all $N \leq 6$ and pairwise-central.} The noise-floor
    argument relies on (i) pairwise central forces and (ii) a symplectic integrator. For
    three-body forces or non-symplectic integrators the momentum and angular floors are
    no longer automatically at roundoff.
  \item \textbf{No accuracy comparison against published numbers.} The HNN baseline was
    trained by us; we did not reproduce a published HNN result. The tie should be read
    as ``our HNN implementation is a fair baseline'', not as a leaderboard claim.
  \item \textbf{The critic is used only for selection.} We do not show that it improves
    \emph{training} (e.g.\ as a reward or regulariser); that is untested.
\end{enumerate}

\section{Retractions}
\label{sec:retractions}

Three claims appeared in earlier internal drafts of this work. All are false, and the
process of falsifying them produced the methodological results of
Sections~\ref{sec:metric}, \ref{sec:ablation} and \ref{sec:taxonomy}.

\paragraph{R1. ``The conservation critic is more robust out of distribution than a
learned critic.''} False. On a family shift the two do not differ (conservation
$+0.001$, learned $-0.005$, overlapping CIs). On a \emph{severity} shift toward subtler
violations the conservation critic degrades \emph{more} ($-11.4$ points versus $-1.0$).
Its advantage is absolute level and zero labelling cost, not distributional robustness.

\paragraph{R2. ``Conservation and learned critics have complementary blind spots;
combining them is consistently better.''} False. An \emph{optimally weighted} fusion
(logistic regression over the critic scores, fitted on held-out training data) assigns
the conservation critic the largest weight ($+1.635$), the supervised error regressor a
\textbf{negative} weight ($-0.612$), and still ties the conservation critic alone
($0.641\times$ versus $0.637\times$ ID; $0.608\times$ versus $0.614\times$ OOD). The
signals are redundant, not complementary. An earlier version of this claim appeared to
hold only because the ``learned'' critic was given conservation residuals as input
features, making the conclusion true by construction.

\paragraph{R3. ``Conservation critics can validate the Lumyn pipeline itself.''} False,
for a structural reason: that pipeline integrates the true equations of motion, so its
output \emph{is} a solution and the failure mode a validator would catch cannot occur. A
real corruption/restoration experiment confirms this --- the validation layer's causal
contribution is exactly zero, because it is never acted upon. A conservation critic is
useful precisely when the generator's output is \emph{not} a solution, which is true of
learned models and false of a physics integrator.

\section{Conclusion}

A critic built from conservation laws, requiring no training data and no labels, ties
the strongest learned critics at selecting good candidates from a learned physical
generator: 100\% within-group accuracy on standard benchmarks (matching a Hamiltonian
Neural Network exactly), and 83--85\% on a 6-body family where a supervised
pairwise-ranking critic trained on 960 labelled candidates reaches 71--72\%.

The reason it works is numerical rather than physical. Conservation channels whose
invariance is a structural property of the pairwise central force --- linear and angular
momentum --- sit at a machine-precision floor and separate true from generated
trajectories by ten or more orders of magnitude. The energy channel, conserved only to
the integrator's truncation error, does not. Choosing channels by their noise floor,
rather than by the physical prominence of the quantity, is what makes the critic work.

We also report what does not work, because it constrains the claim: the critic is not
more robust out of distribution, combining it with learned critics does not help, and
synthetic failure taxonomies are not structurally comparable to real generator failures
unless calibrated first.

\begin{thebibliography}{9}

\bibitem{chenciner2000}
A.~Chenciner and R.~Montgomery.
\newblock A remarkable periodic solution of the three-body problem in the case of equal
masses.
\newblock \emph{Annals of Mathematics}, 152(3):881--901, 2000.

\bibitem{cranmer2020lNN}
M.~Cranmer, S.~Greydanus, S.~Hoyer, P.~Battaglia, D.~Spergel, and S.~Ho.
\newblock Lagrangian Neural Networks.
\newblock \emph{arXiv:2003.04630}, 2020.

\bibitem{greydanus2019hnn}
S.~Greydanus, M.~Dzamba, and J.~Yosinski.
\newblock Hamiltonian Neural Networks.
\newblock \emph{Advances in Neural Information Processing Systems}, 2019.

\end{thebibliography}

\appendix

\section{Reproducibility}

\begin{table}[h]
\centering
\caption{Result-to-script map.}
\label{tab:repro}
\begin{tabular}{ll}
\toprule
Result & Script \\
\midrule
Selection metrics, 6-body, 5 seeds & \texttt{experiments/experiment\_h\_benchmark.py} \\
Channel ablation & \texttt{experiments/experiment\_e\_channel\_combos.py} \\
Standard benchmarks $+$ HNN & \texttt{experiments/experiment\_k\_hnn\_equation\_critic.py} \\
Strong supervised baselines & \texttt{experiments/experiment\_c\_strong\_baselines.py} \\
Failure-mode comparability & \texttt{experiments/experiment\_i\_failure\_modes.py} \\
Noise-floor sweep & \texttt{experiments/experiment\_g\_error\_functionals.py} \\
HNN fairness diagnostics & \texttt{experiments/diagnose\_hnn\_*.py} \\
Channel-set regression tests & \texttt{tests/test\_conservation\_critic.py} (18 tests) \\
Checker-defect regression tests & \texttt{tests/test\_conservation\_checker.py} (11 tests) \\
\bottomrule
\end{tabular}
\end{table}

Test suite at the time of writing: \textbf{108 tests, 0 failures} (1 optional-dependency
skip).

\end{document}
